\documentclass[12pt, titlepage]{article}

\usepackage{booktabs}
\usepackage{tabularx}
\usepackage{hyperref}
\hypersetup{
    colorlinks,
    citecolor=black,
    filecolor=black,
    linkcolor=red,
    urlcolor=blue
}
\usepackage[round]{natbib}

\input{../Comments}
\input{../Common}

\begin{document}

\title{Verification and Validation Report: \progname} 
\author{\authname}
\date{\today}
	
\maketitle

\pagenumbering{roman}

\section{Revision History}

\begin{tabularx}{\textwidth}{p{3cm}p{2cm}X}
\toprule {\bf Date} & {\bf Version} & {\bf Notes}\\
\midrule
Date 1 & 1.0 & Notes\\
Date 2 & 1.1 & Notes\\
\bottomrule
\end{tabularx}

~\newpage

\section{Symbols, Abbreviations and Acronyms}

\renewcommand{\arraystretch}{1.2}
\begin{tabular}{l l} 
  \toprule		
  \textbf{symbol} & \textbf{description}\\
  \midrule 
  T & Test\\
  \bottomrule
\end{tabular}\\

\wss{symbols, abbreviations or acronyms -- you can reference the SRS tables if needed}

\newpage

\tableofcontents

\listoftables %if appropriate

\listoffigures %if appropriate

\newpage

\pagenumbering{arabic}

This document ...

\section{Functional Requirements Evaluation}

\textit{Note: While the original VnV Plan specified automated Node console harnesses for Functional Requirements testing (e.g., ST-GM-01), the final evaluation was conducted via \textbf{Manual UI-based Testing} on the deployed application. This transition was made to better validate the end-to-end integration of the frontend, backend, and real-time WebSocket engine under realistic network conditions.}

\subsection{Game Manager: Start, Turn, Rules, Specials, End (FR-1..FR-5)}

\paragraph{Start new game initializes state (ST-GM-01)}
\begin{itemize}
  \item \textbf{Requirement:} FR-1
  \item \textbf{Test Type:} Manual UI
  \item \textbf{Expected Output:} UI transitions to the Game Screen, both avatars show 7 cards in hand, discard pile shows 1 top card, and UI visually indicates it is User A's turn.
  \item \textbf{Actual Output:} \wss{Fill in your actual observation here (e.g., "Observed 7 cards per player, UI highlights Player A's turn, top card is visible.")}
  \item \textbf{Result:} \wss{PASS / FAIL}
  \item \textbf{Notes:} \wss{Any deviations or extra notes}
\end{itemize}

\paragraph{Legal move: same suit (ST-GM-02)}
\begin{itemize}
  \item \textbf{Requirement:} FR-3
  \item \textbf{Test Type:} Manual UI
  \item \textbf{Expected Output:} Card of the same suit moves to discard pile and turn indicator shifts to User B.
  \item \textbf{Actual Output:} \wss{Fill in actual observation}
  \item \textbf{Result:} \wss{PASS / FAIL}
\end{itemize}

\paragraph{Legal move: sum equals 12 base-12 (ST-GM-03)}
\begin{itemize}
  \item \textbf{Requirement:} FR-3
  \item \textbf{Test Type:} Manual UI
  \item \textbf{Expected Output:} Card whose decimal value adds up to 12 with the top card is accepted and turn advances.
  \item \textbf{Actual Output:} \wss{Fill in actual observation}
  \item \textbf{Result:} \wss{PASS / FAIL}
\end{itemize}

\paragraph{Special card: wild ten suit selection (ST-GM-04)}
\begin{itemize}
  \item \textbf{Requirement:} FR-4
  \item \textbf{Test Type:} Manual UI
  \item \textbf{Expected Output:} Selecting a wildcard prompts for a new suit. UI updates to show forced suit, restricting next player.
  \item \textbf{Actual Output:} \wss{Fill in actual observation}
  \item \textbf{Result:} \wss{PASS / FAIL}
  \item \textbf{Notes:} Note that implementation uses rank "10" as wild instead of "8" as originally written in VnVPlan.
\end{itemize}

\paragraph{Illegal move is rejected with guidance (ST-GM-05)}
\begin{itemize}
  \item \textbf{Requirement:} FR-3
  \item \textbf{Test Type:} Manual UI
  \item \textbf{Expected Output:} Illegal cards are not highlighted and cannot be selected or played.
  \item \textbf{Actual Output:} \wss{Fill in actual observation}
  \item \textbf{Result:} \wss{PASS / FAIL}
\end{itemize}

\paragraph{End of game detection (ST-GM-06)}
\begin{itemize}
  \item \textbf{Requirement:} FR-5
  \item \textbf{Test Type:} Manual UI
  \item \textbf{Expected Output:} Successfully playing final card transitions UI to a summary screen declaring the winner and displaying total scores.
  \item \textbf{Actual Output:} \wss{Fill in actual observation}
  \item \textbf{Result:} \wss{PASS / FAIL}
\end{itemize}


\subsection{Score Features: Dual-base, Highlights (FR-6..FR-9)}

\paragraph{Score calculated and shown in decimal and Dozenal (ST-SC-01)}
\begin{itemize}
  \item \textbf{Requirement:} FR-6
  \item \textbf{Test Type:} Manual UI
  \item \textbf{Expected Output:} Round Result popup correctly summarizes lost hand points in both Decimal and Dozenal (e.g., 10$_{12}$) notations side by side.
  \item \textbf{Actual Output:} \wss{Fill in actual observation}
  \item \textbf{Result:} \wss{PASS / FAIL}
\end{itemize}

\paragraph{Valid moves highlighted exactly (ST-SC-03)}
\begin{itemize}
  \item \textbf{Requirement:} FR-9
  \item \textbf{Test Type:} Manual UI
  \item \textbf{Expected Output:} Exactly the cards that are legal to play are visually highlighted/raised, while others are dimmed/disabled.
  \item \textbf{Actual Output:} \wss{Fill in actual observation}
  \item \textbf{Result:} \wss{PASS / FAIL}
\end{itemize}

\paragraph{Primary base toggle affects labeling only (ST-SC-04)}
\begin{itemize}
  \item \textbf{Requirement:} FR-7
  \item \textbf{Test Type:} Manual UI
  \item \textbf{Expected Output:} Score displays correctly reflect dual-base formats consistently throughout the UI.
  \item \textbf{Actual Output:} \wss{Fill in actual observation}
  \item \textbf{Result:} \wss{PASS / FAIL}
\end{itemize}


\subsection{Login Manager: Account, Sessions, Guest, Validation (FR-10..FR-13)}

\paragraph{Account creation creates unique user and session (ST-LM-01)}
\begin{itemize}
  \item \textbf{Requirement:} FR-10
  \item \textbf{Test Type:} Manual UI
  \item \textbf{Expected Output:} Valid registration logs user in. Duplicate registration is rejected with "username taken" error.
  \item \textbf{Actual Output:} \wss{Fill in actual observation}
  \item \textbf{Result:} \wss{PASS / FAIL}
\end{itemize}

\paragraph{Login and logout preserve progress (ST-LM-02)}
\begin{itemize}
  \item \textbf{Requirement:} FR-11
  \item \textbf{Test Type:} Manual UI
  \item \textbf{Expected Output:} Logging out and logging back in retains exact match history in Profile page.
  \item \textbf{Actual Output:} \wss{Fill in actual observation}
  \item \textbf{Result:} \wss{PASS / FAIL}
\end{itemize}

\paragraph{Guest mode creates temporary session (ST-LM-03)}
\begin{itemize}
  \item \textbf{Requirement:} FR-12
  \item \textbf{Test Type:} Manual UI
  \item \textbf{Expected Output:} "Play as Guest" creates temporary profile. Logging out deletes profile, preventing future logins to that guest account.
  \item \textbf{Actual Output:} \wss{Fill in actual observation}
  \item \textbf{Result:} \wss{PASS / FAIL}
\end{itemize}

\paragraph{Credential validation gates access (ST-LM-04)}
\begin{itemize}
  \item \textbf{Requirement:} FR-13
  \item \textbf{Test Type:} Manual UI
  \item \textbf{Expected Output:} Incorrect password rejected with "invalid credentials". Correct password proceeds to Lobby.
  \item \textbf{Actual Output:} \wss{Fill in actual observation}
  \item \textbf{Result:} \wss{PASS / FAIL}
\end{itemize}


\subsection{Data Manager: Storage, Retrieval, Update, Deletion (FR-14..FR-17)}

\paragraph{User data is stored on events that require persistence (ST-DM-01)}
\begin{itemize}
  \item \textbf{Requirement:} FR-14
  \item \textbf{Test Type:} Manual UI
  \item \textbf{Expected Output:} Completing a game adds a correct entry to the Profile Match History.
  \item \textbf{Actual Output:} \wss{Fill in actual observation}
  \item \textbf{Result:} \wss{PASS / FAIL}
\end{itemize}

\paragraph{Stored data can be retrieved on login or profile view (ST-DM-02)}
\begin{itemize}
  \item \textbf{Requirement:} FR-15
  \item \textbf{Test Type:} Manual UI
  \item \textbf{Expected Output:} Profile page successfully fetches and renders stats (win rate) and match history.
  \item \textbf{Actual Output:} \wss{Fill in actual observation}
  \item \textbf{Result:} \wss{PASS / FAIL}
\end{itemize}

\paragraph{Data updates after games and profile changes (ST-DM-03)}
\begin{itemize}
  \item \textbf{Requirement:} FR-16
  \item \textbf{Test Type:} Manual UI
  \item \textbf{Expected Output:} Changing Display Name in settings persists across full page refreshes.
  \item \textbf{Actual Output:} \wss{Fill in actual observation}
  \item \textbf{Result:} \wss{PASS / FAIL}
\end{itemize}

\paragraph{User can permanently delete data (ST-DM-04)}
\begin{itemize}
  \item \textbf{Requirement:} FR-17
  \item \textbf{Test Type:} Manual UI
  \item \textbf{Expected Output:} Clicking "Delete Account" forces logout and permanently prevents logging back in with old credentials.
  \item \textbf{Actual Output:} \wss{Fill in actual observation}
  \item \textbf{Result:} \wss{PASS / FAIL}
\end{itemize}

\section{Nonfunctional Requirements Evaluation}

\subsection{Usability}
		
\subsection{Performance}

\subsection{etc.}
	
\section{Comparison to Existing Implementation}	

This section will not be appropriate for every project.

\section{Unit Testing}

\section{Changes Due to Testing}

\wss{This section should highlight how feedback from the users and from 
the supervisor (when one exists) shaped the final product.  In particular 
the feedback from the Rev 0 demo to the supervisor (or to potential users) 
should be highlighted.}

\section{Automated Testing}
		
\section{Trace to Requirements}
		
\section{Trace to Modules}		

\section{Code Coverage Metrics}

\bibliographystyle{plainnat}
\bibliography{../../refs/References}

\newpage{}
\section*{Appendix --- Reflection}

The information in this section will be used to evaluate the team members on the
graduate attribute of Reflection.

\input{../Reflection.tex}

\begin{enumerate}
  \item What went well while writing this deliverable? 
  \item What pain points did you experience during this deliverable, and how
    did you resolve them?
  \item Which parts of this document stemmed from speaking to your client(s) or
  a proxy (e.g. your peers)? Which ones were not, and why?
  \item In what ways was the Verification and Validation (VnV) Plan different
  from the activities that were actually conducted for VnV?  If there were
  differences, what changes required the modification in the plan?  Why did
  these changes occur?  Would you be able to anticipate these changes in future
  projects?  If there weren't any differences, how was your team able to clearly
  predict a feasible amount of effort and the right tasks needed to build the
  evidence that demonstrates the required quality?  (It is expected that most
  teams will have had to deviate from their original VnV Plan.)
\end{enumerate}

\end{document}